\documentclass[11pt,letterpaper]{article}
\usepackage[T1]{fontenc}
\usepackage[utf8]{inputenc}
\usepackage{lmodern,microtype}
\usepackage[margin=0.78in,headheight=15pt]{geometry}
\usepackage{amsmath,amssymb,booktabs,longtable,tabularx,array,enumitem,xcolor,hyperref,fancyhdr,lastpage,needspace}
\definecolor{EngineBlue}{HTML}{183B56}
\definecolor{EngineTeal}{HTML}{117A83}
\definecolor{EngineGray}{HTML}{596671}
\definecolor{EngineLight}{HTML}{EEF4F7}
\hypersetup{colorlinks=true,urlcolor=EngineTeal,linkcolor=EngineBlue}
\setlength{\parindent}{0pt}
\setlength{\parskip}{5pt}
\emergencystretch=3em
\setlist{nosep,leftmargin=*}
\pagestyle{fancy}\fancyhf{}
\fancyhead[L]{\small\sffamily Game Engine Architecture | Engineering Suite}
\fancyhead[R]{\small\sffamily \DocType}
\fancyfoot[L]{\scriptsize Source-grounded engineering curriculum | v1.0}
\fancyfoot[R]{\scriptsize \thepage/\pageref{LastPage}}
\newcommand{\DocType}{Reference}
\newcommand{\SuiteHeading}[3]{%
\renewcommand{\DocType}{#2}%
{\small\bfseries\sffamily\color{EngineTeal}#2\hfill #3}\par\vspace{6pt}
{\LARGE\bfseries\color{EngineBlue}#1}\par\vspace{4pt}
{\color{EngineTeal}\rule{\linewidth}{1pt}}\par\vspace{7pt}}
\newcommand{\LabEvidence}[1]{\par\textbf{Evidence required:} #1\par}
\newcommand{\WriteLines}[1]{\par\begingroup\setlength{\parskip}{9pt}\count255=0\loop\ifnum\count255<#1\noindent\rule{\linewidth}{0.35pt}\par\advance\count255 by1\repeat\endgroup\smallskip}
\newcommand{\WritingBox}[1]{\noindent\fbox{\parbox[t][#1][t]{0.97\linewidth}{\mbox{}}}\par\smallskip}
\newcommand{\SourceNote}{\footnotesize\textit{Primary reading: Jason Gregory, \emph{Game Engine Architecture}, 3rd ed. (CRC Press, 2018). The narrative, exercises, laboratory designs and scoring instruments in this suite are independently authored teaching materials, not a reproduction of the textbook.}\normalsize\par}
\newcommand{\Criterion}[2]{\textbf{#1.} #2\par}
\newcommand{\FormField}[2]{\Needspace{3\baselineskip}\textbf{#1}\par\textcolor{EngineGray}{\footnotesize #2}\par\WriteLines{2}}
\newcommand{\BlankField}[1]{\Needspace{3\baselineskip}\textbf{#1}\par\WriteLines{2}}
\newcolumntype{Y}{>{\raggedright\arraybackslash}X}

\begin{document}
\SuiteHeading{Physics and Audio Integration Workbook — Guided Template}{Form / Template}{Controlled form}
\section{Purpose and use}
This is the guided version of the form. For each numbered field, use the accompanying prompt to capture evidence. Leave no required criterion unsubstantiated; cite test IDs and measured data rather than subjective claims.

\subsection{Scenario}
Guidance: World units, gravity, bodies, contact materials and listener location.

\WritingBox{1.05in}
\subsection{Physics steps}
Guidance: dt, substeps, collision detection, restitution and ownership.

\WritingBox{1.05in}
\subsection{Audio configuration}
Guidance: Sample rate, channels, buffer size, gain and attenuation.

\WritingBox{1.05in}
\subsection{Event lifecycle}
Guidance: Contact to gameplay event to audio command to output buffer.

\WritingBox{1.05in}
\subsection{Stability and continuity}
Guidance: Penetration, energy drift, clipping, underruns and callback cost.

\WritingBox{1.05in}
\subsection{Failure injection}
Guidance: Invalid body, missing sound, delayed asset and bounded queue saturation.

\WritingBox{1.05in}
\subsection{Conclusion}
Guidance: Correctness, performance and outstanding integration risks.

\WritingBox{1.05in}
\section{Review record}
\BlankField{Prepared by / date}
\BlankField{Reviewer / decision / date}

\end{document}
